% Pursuing Stacks (cotes 134-2 to 134-8): archive cote 134-6, site page 5 (PDF page 6), typescript page 324.
% typescript. Draft compared: ps5.tex:1571-1644.
% AI-transcribed by Claude Opus 5.5 (claude-opus-5-5, adaptive thinking, effort high), Message Batches API, batch ps3 (msgbatch_01XdFmuVVBAcgc7EFxdtCqsj), 2026-11-12; not yet checked by eye.
\scanpage{134-6/5}
\mspage{324}

A  is either \hand{$\square$} , or \hand{$\bigcirc$} .  In these three cases,  the functor (6) can
be  factorized in a natural way through \deleted{\unsure{a}}\inserted{the} category
\[ \mathrm{Hotab}_A \;=\; W^{-1}A^{\wedge}_{ab} \qquad , \]
where now W stands for the set of ``weak equivalences'' in  $A^{\wedge}_{ab}$ ,  defined
in the same way as above in the case $A=$\hand{$\Delta$}.  The question then arises,
for any given A,  as to whether such a factorization can be still obtai-%
ned,  and how exactly.

This formulation is inspired by the description of abelianization
of homotopy types via the (slightly sophisticated) diagram (3).  When
following the more naïve approach of the remark above,  this leads us to
the closely related question of defining (6)  via a composition
\[ (9) \qquad A^{\wedge} \xrightarrow{\ \mathrm{Wh}_A\ } A^{\wedge}_{ab} \xrightarrow{\ \hand{$L$}\ } \mathrm{Ch.(Ab)} \qquad , \]
(for suitable functor \hand{$L$} , cf.below),  by passing to localizations.

Both approaches seem to me of interest.  The first one,  to make a
sense at all as stated,  relies on the existence of a canonical functor
\[ (9') \qquad \mathrm{Hot}_A \longrightarrow \mathrm{Hotab}_A \qquad , \]
induced by the abelianization functor
\[ \mathrm{Wh}_A \ : \quad A^{\wedge} \longrightarrow A^{\wedge}_{ab} \quad ,X \longmapsto \underline{Z}^{(X)} \quad , \]
i.e.  on the \uline{validity of Whitehead's theorem},  \uline{with \hand{$\Delta$} replaced by} A.
This looks like an interesting question,  whose answer should be in the
affirmative.  At any rate,  if we can find an aspheric functor
\[ (10) \qquad j: \hand{$\Delta$} \longrightarrow A^{\wedge} \]
(with respect to the standard asphericity structure of A )\xout{maybe we bet-}
\xout{ter \inserted[typed]{A} assume aspheric then}  ,  then the answer is affirmative,  as we are
immediately reduced to the known case when A is replaced by \hand{$\Delta$}.  Thus,
the answer is possibly tied up with the question of existence of test
functors,  which we'll deal with presumably in part V.   It should be \xout{xxx}
noted though that if such a functor (10) exists,  then necessarily A is
aspheric,  and even totally aspheric  -  a substantial restriction indeed\unsure{!}

M\over{e}{o}re generally,  let
\[ (10') \qquad j^{\unsure{*}}: \ B \longrightarrow A^{\wedge} \]
any aspheric functor with respect to the standard asphericity structure
of $A^{\wedge}$, where B is any small category.   \xout{Wxxk}  The corresponding functor

\margin{5}
